
%%%%%%%%%%  scan idx 12  |  folio 1  %%%%%%%%%%


SGA 7


\underline{Exposé I}


\underline{\underline{Resumen de las primeras exposiciones de A. Grothendieck}}


\underline{redactado por P. Deligne}


\underline{Sumario}

0. Preliminares

1. Demostración aritmética del teorema de monodromía

2. Ciclos evanescentes

3. Demostración geométrica del teorema de monodromía

4. Criterios de anulación para los haces de ciclos evanescentes

5. Acción de la monodromía sobre los $\pi_1$

6. Apéndice de P. Deligne: demostración aritmética del
teorema de reducción estable


0. \underline{Preliminares}

0.0. \underline{Terminología}. Recordemos las definiciones siguientes:

(0.0.1)    \underline{rasgo}: espectro de un anillo de valoración discreta. Los puntos genérico y cerrado de un rasgo se denotan a menudo por $\eta$ y $s$.

(0.0.2)    \underline{estrictamente local}: para un anillo: local henseliano con cuerpo residual separablemente cerrado. Para un esquema: espectro de un anillo de esta clase.

(0.0.3)    \underline{punto geométrico de} $S$ (en esta exposición): morfismo $\bar{x} : \mathrm{Spec}(\bar{k}) \longrightarrow S$ de imagen $x \in S$, siendo $k(\bar{x}) \underset{\mathrm{dfn}}{=} \bar{k}$ una clausura separable de $k(x)$. Para un rasgo henseliano $S = \mathrm{Spec}(V)$, se denota a menudo por $\bar{\eta}$ un punto geométrico genérico (es decir, de imagen $\eta$) de $S$. El cuerpo residual del anillo de valoración $\bar{V}$, clausura entera de $V$ en $k(\bar{\eta})$, es una extensión algebraica separablemente cerrada de $k(s)$ y define un punto geométrico $\bar{s}$ de $S$ localizado en $s$.


%%%%%%%%%%  scan idx 13  |  folio 2  %%%%%%%%%%

(0.0.4)    \underline{localización estricta de} $X$ \underline{en el punto geométrico} $\bar{x}$: SGA 4 VIII 4.4.

(0.0.5)    \underline{esencialmente de tipo finito sobre} $S$: para un $S$-esquema local (resp. local henseliano, resp. estrictamente local): espectro de un anillo local (resp. del henselizado de un anillo local, resp. localización estricta) de un esquema de tipo finito sobre $S$.

(0.0.6)    $F \longmapsto F(n)$: giro de Tate.

(0.0.7)    Un endomorfismo $T$ de un espacio vectorial de dimensión finita es \underline{cuasiunipotente} si sus valores propios son raíces de la unidad, es decir, si existen $N \geqslant 1$, $M \geqslant 1$ tales que $(T^N - 1)^M = 0$.


0.1.  \underline{Recordatorio sobre la cohomología $\ell$-ádica}

Sea $X$ un esquema de tipo finito sobre un cuerpo algebraicamente cerrado $k$ de exponente característico $p$, y sea $\ell$ un número primo distinto de $p$.

Los grupos de cohomología $\ell$-ádica de $X$ fueron definidos en SGA 4 y SGA 5.

a) $H^i(X,\mathbb{Z}/\ell^n)$ es el $i^{\text{ésimo}}$ grupo de cohomología del sitio étale $X_{et}$ de $X$ con valores en $\mathbb{Z}/\ell^n$;

b) por definición, $H^i(X,\mathbb{Z}_\ell) = \underset{n}{\underleftarrow{\ell\mathrm{im}}}\, H^i(X,\mathbb{Z}/\ell^n)$;

c) por definición, $H^i(X,\mathbb{Q}_\ell) = H^i(X,\mathbb{Z}_\ell) \otimes_{\mathbb{Z}_\ell} \mathbb{Q}_\ell$.

Los grupos de cohomología $\ell$-ádica con soportes propios de $X$ se definen del mismo modo a partir de los $H^i_c(X,\mathbb{Z}/\ell^n)$ de SGA 4 XVII.

En cada uno de los casos siguientes, se sabe demostrar que estos grupos $H^i(X)$ poseen propiedades de finitud razonables:

(0.1.1)    en cohomología con soporte propio;

(0.1.2)    para $X$ propio sobre $k$;

(0.1.3)    para $X$ liso sobre $k$ (gracias a a) y a la dualidad de Poincaré, SGA 4 XVIII);

(0.1.4)    cuando se dispone de la resolución de singularidades al modo de Hironaka para los esquemas de tipo finito sobre $k$, de dimensión $\leqslant \dim(X)$;

(0.1.5)    para $i = 0, 1$.


%%%%%%%%%%  scan idx 14  |  folio 3  %%%%%%%%%%

En todos estos casos, los $H^i(X,\mathbb{Z}/\ell^n)$ son finitos, los $H^i(X,\mathbb{Z}_\ell)$ son de tipo finito, los $H^i(X,\mathbb{Q}_\ell)$ son de dimensión finita y (salvo quizá para (0.1.5)) se dispone de sucesiones exactas cortas escindidas

\begin{equation*}
0 \longrightarrow H^i(X,\mathbb{Z}_\ell) \otimes \mathbb{Z}/\ell^n \longrightarrow H^i(X,\mathbb{Z}/\ell^n) \longrightarrow \mathrm{Tor}_1(H^{i+1}(X,\mathbb{Z}_\ell),\mathbb{Z}/\ell^n) \longrightarrow 0 \ .
\end{equation*}
Además, la demostración del teorema de finitud proporciona en cada caso una demostración de un teorema de constructibilidad genérica: si $k$ es el cuerpo de fracciones de un esquema noetheriano entero $S$ y $X$ es la fibra genérica de $f_1 : X_1 \longrightarrow S$ (de tipo finito), existe un abierto no vacío $U \subset S$ tal que, sobre $U$, los $R^i f_* \mathbb{Z}/\ell^n$ sean localmente constantes y su formación sea compatible con todo cambio de base.

0.2. Dejemos de suponer que $k$ es algebraicamente cerrado, y sea $\bar{k}$ una clausura algebraica (o separable, lo que aquí sería equivalente) de $k$. Consideraremos los grupos de cohomología «geométricos» $H^i(X_{\bar{k}})$ ($X_{\bar{k}}$ es el esquema deducido de $X$ por extensión de escalares de $k$ a $\bar{k}$). Por transporte de estructura, el grupo de Galois $\mathrm{Gal}(\bar{k}/k)$ actúa sobre estos grupos. En cohomología con coeficientes en $\mathbb{Q}_\ell$, cuando se está en uno de los casos (0.1.1) a (0.1.5), se obtienen así representaciones continuas de $\mathrm{Gal}(\bar{k}/k)$ en un grupo lineal $\mathrm{GL}(n,\mathbb{Q}_\ell)$.


0.3. \underline{Recordatorio sobre los rasgos} (0.0.1)

Para las demostraciones, remitimos a $[4]$. Sean $S$ un rasgo henseliano, y $\eta$, $s$, $V$, $\bar{V}$, $\bar{\eta}$, $\bar{s}$ como en (0.0.1) y (0.0.3). Denotemos por $p$ el exponente característico de $k(s)$ y por $\mathrm{Gal}(\bar{\eta}/\eta)$ (resp. $\mathrm{Gal}(\bar{s}/s)$) el grupo de Galois $\mathrm{Gal}(k(\bar{\eta})/k(\eta))$ (resp. ...). Por transporte de estructura, $\mathrm{Gal}(\bar{\eta}/\eta)$ actúa sobre $k(\bar{s})$, de donde se obtiene una aplicación de $\mathrm{Gal}(\bar{\eta}/\eta)$ en (de hecho, sobre) $\mathrm{Gal}(\bar{s}/s)$, cuyo núcleo es el \underline{grupo de inercia} $I$. El subcuerpo de $k(\bar{\eta})$ definido por $I$ es una extensión no ramificada máxima de $k(\eta)$.

El grupo profinito $I$ admite un mayor subgrupo $P$ que es un pro-$p$-grupo. El cociente $I/P$ es el \underline{grupo de inercia moderada}. Canónicamente se tiene

\begin{equation*}
I/P \simeq \underset{n}{\underleftarrow{\lim}}\ \mu_n(k(\bar{s})) \underset{\mathrm{dfn}}{=} \widehat{\mathbb{Z}}(1)(k(\bar{s})) \ .
\end{equation*}

%%%%%%%%%%  scan idx 15  |  folio 4  %%%%%%%%%%

Si $\phi_n$ es la aplicación canónica de $I/P$ en $\mu_n(\bar{k})$, y $u$ es un elemento de $\bar{V}$ de valoración $1/n$, se tiene, para $\sigma \in I$,

\begin{equation*}
\sigma(u) \;\equiv\; \phi_n(\sisi{u}{\sigma}).u \qquad (\text{mód. elementos de valoración} > 1/n).
\end{equation*}
En resumen, se tiene:

\begin{equation}\tag{0.3.1}
\begin{aligned}
\operatorname{Gal}(\bar\eta/\eta) &\supset I \supset P \ , \\
\operatorname{Gal}(\bar\eta/\eta)/I &\simeq \operatorname{Gal}(\bar s/s) \ , \\
I/P &\simeq \hat{\mathbf{Z}}(1)(k(\bar s)) \ , \\
P &\quad : \text{grupo pro-p} \ (\ \{e\} \ \text{si} \ p = 1 \ ) \ ,
\end{aligned}
\end{equation}
y la acción (por automorfismos interiores en $\operatorname{Gal}(\bar\eta/\eta)$) de $\operatorname{Gal}(\bar\eta/\eta)/I$ sobre $I/P$ se identifica con la acción natural de $\operatorname{Gal}(\bar s/s)$ sobre $\hat{\mathbf{Z}}(1)(k(\bar s))$.

0.4.  Sea, más generalmente, $S$ estrictamente local y regular, $D$ un divisor regular en $S$ y $p'$ el exponente característico del punto genérico de $D$. Sean $\sisi{\eta}{\bar\eta}$ un punto geométrico genérico de $S$ y $\bar s$ el punto geométrico localizado en el punto cerrado $s$ que se deduce de él. Según el lema de Abhyankar, el grupo fundamental $\pi_1(S\!-\!D,\bar\eta)$ admite una descomposición análoga a (0.3.1):

\begin{equation}\tag{0.4.1}
\begin{aligned}
\pi_1(S\!-\!D,\bar\eta) &\supset I \supset P \\
\pi_1(S\!-\!D,\bar\eta)/I &\simeq \operatorname{Gal}(\bar s/s) \\
I/P &\simeq \hat{\mathbf{Z}}(1)(k(\bar s)) \\
P &\quad : \text{sin cociente de orden primo con } p' \ .
\end{aligned}
\end{equation}
La acción por automorfismos interiores de $\pi_1(S\!-\!D,\bar\eta)/I$ sobre $I/P$ se identifica con la acción natural de $\operatorname{Gal}(\bar s/s)$ sobre $\hat{\mathbf{Z}}(1)(k(\bar s))$.


0.5.  \underline{La desingularización de Néron}

El método de desingularización de Néron proporciona el resultado siguiente:

\underline{Lema} 0.5.1.  \underline{Sea} $S \longrightarrow S_o$ \underline{un morfismo de rasgos henselianos,} \underline{con} $S = Spec(V)$, $S_o = Spec(V_o)$ \underline{y} $\pi$ \underline{un uniformizante de} $V_o$. \underline{Supongamos que} $\pi$ \underline{es un uniformizante de} $V$ \underline{y que las extensiones} $k(s)/k(s_o)$ \underline{y} $k(\eta)/k(\eta_o)$ \underline{son separables.} \underline{Entonces,} $V$ \underline{es límite inductivo de} $V_o$\underline{-álgebras henselianas lisas y esencialmente de tipo finito} $B_i$ \underline{sobre} $V_o$.


%%%%%%%%%%  scan idx 16  |  folio 5  %%%%%%%%%%

Para construir las $V_o$-álgebras buscadas, se toma $V_o \subset B \subset V$, con $B$ de tipo finito, y se desingulariza $Spec(B)$ al modo de Néron ([1] § 4).

Sea $S = Spec(V)$ un rasgo henseliano.

(0.5.2)  Si $S$ es de igual característica, con uniformizante $\pi$, se puede aplicar (0.5.1) para $S_o = \mathbf{F}_p[\pi]_{(\pi)}$. La extensión $s/s_o$ es separable porque $\mathbf{F}_p$ es perfecto; la extensión $\eta/\eta_o$ lo es porque $\pi$, al ser un uniformizante, no es una potencia $p^{\text{ésima}}$ (si $p \neq 1$) y, por tanto, forma parte de una $p$-base. Se obtiene que $S$ es límite de espectros de $\mathbf{F}_p[\pi]_{(\pi)}$-álgebras henselianas lisas esencialmente de tipo finito.

(0.5.3)  Si $S$ es completo, de característica mixta y con cuerpo residual perfecto $k$, $V$ es una extensión de Eisenstein del anillo de vectores de Witt ($=$ anillo de Cohen) $W(k)$

\begin{equation*}
V \simeq W(k) \ [\pi] \ /(\pi^n + \sum_0^{\,n-1} a_i \ \pi^i) \qquad .
\end{equation*}
Si se perturban un poco los $a_i$, se obtiene una extensión isomorfa. Podemos, pues, suponer que los $a_i$ pertenecen a $W(k')$, con $k'$ de tipo finito sobre $\mathbf{F}_p$. Sea $k_o$ la clausura perfecta de $k'$ en $k$ y

\begin{equation*}
V_o = W(k_o) \ [\pi]/(\pi^n + \sum_0^{\,n-1} a_i \ \pi^i) \qquad .
\end{equation*}
El lema 0.5.1 se aplica a $V/V_o$, y el cuerpo residual de $V_o$ es una extensión puramente inseparable de un cuerpo de tipo finito sobre $\mathbf{F}_p$.


1.  \underline{Demostración aritmética del teorema de monodromía}

El lector encontrará en [5], Apéndice, una demostración del siguiente resultado de A. Grothendieck.

\underline{Proposición} 1.1.  \underline{Sean} $S$ \underline{un rasgo henseliano,} $\ell$ \underline{un número primo distinto del exponente característico} $p$ \underline{de} $k(s)$ \underline{y} $\rho$ \underline{una representación continua de} $\operatorname{Gal}(\bar\eta/\eta)$ \underline{en} $GL(n,\mathbf{Q}_\ell)$. \underline{Supongamos que se cumple la condición siguiente}:

$(*_\ell)$      \underline{Ninguna extensión finita de} $k(s)$ \underline{contiene todas las raíces de la unidad cuyo orden es una potencia de} $\ell$.


%%%%%%%%%%  scan idx 17  |  folio 6  %%%%%%%%%%

\underline{Entonces existe un subgrupo abierto} $I_1$ \underline{del grupo de inercia} $I$ \underline{tal que} $\rho(\sigma)$ \underline{sea unipotente para} $\sigma \in I_1$.

La condición $(*_\ell)$ se cumple automáticamente para $k(s)$ de tipo finito sobre el cuerpo primo (o puramente inseparable sobre un cuerpo de esta clase).

El teorema siguiente se deduce inmediatamente de 1.1.

\underline{Teorema de monodromía} 1.2.  \underline{Con las notaciones precedentes,} \underline{sea} $X$ \underline{un esquema de tipo finito sobre} $\eta$ \underline{y} $H$ \underline{un espacio de cohomología de} $X_{\bar\eta}$, con \underline{coeficientes en} $\mathbf{Q}_\ell$, \underline{de uno de los tipos} \sisi{(1.1.1) \underline{a} (1.1.5)}{(0.1.1) \underline{a} (0.1.5)}. \underline{Si la condición} $(*_\ell)$ \underline{se cumple,} \underline{la acción de cualquier elemento de} $I$ \underline{sobre} $H$ \underline{es cuasiunipotente.}

Los resultados de paso al límite de 0.5 permiten prescindir de la hipótesis $(*_\ell)$.

\underline{Variante} 1.3.  \underline{La condición} $(*_\ell)$ \underline{de} 1.2 \underline{es superflua.}

Sea $H'_o$ un espacio de $\mathbf{Z}_\ell$-cohomología de $X_{\bar\eta}$ que da lugar a $H$, y $H_o = H'_o$/torsión. Por hipótesis se tiene $H = H_o \otimes_{\mathbf{Z}_\ell} \mathbf{Q}_\ell$, y la acción de $\operatorname{Gal}(\bar\eta/\eta)$ se deduce de

\begin{equation*}
\rho \ : \ \operatorname{Gal}(\bar\eta/\eta) \longrightarrow GL(H_o) \qquad .
\end{equation*}
Sea $I_{\ell'}$ el mayor subgrupo primo con $\ell$ de $I$. Su imagen en $GL(H_o)$ es un grupo de Lie $\ell$-ádico de orden primo con $\ell$ y, por tanto, es finita.

Procedamos a una extensión de rasgos $S' \longrightarrow S$. Esto no cambia $H_o$ (SGA 4 XVI 6.6), y $\rho(I')$ es de índice finito en $\rho(I)$, pues $I'/P'$ es de índice finito en $I/P$ y por lo anterior. Basta, pues, demostrar el teorema sobre $S'$: podemos suponer que

a) $\operatorname{Gal}(\bar\eta/\eta)$ actúa trivialmente sobre $H_o/\ell H_o$;

b) si $S$ es de característica mixta, $S$ es completo con cuerpo residual perfecto (EGA $\mathrm{O}_{III}$ 10.3.1).

Pongamos $S = Spec(V)$, definamos $V_o$ como en (0.5.2) y (0.5.3), y apliquemos (0.5.1). Se tiene

\begin{equation*}
V = \varinjlim B_i
\end{equation*}

%%%%%%%%%%  scan idx 18  |  folio 7  %%%%%%%%%%

con $S_i = Spec(B_i)$ esencialmente liso y de tipo finito sobre $V_o$. Sea $D_i \subset S_i$ de ecuación $\pi = 0$. Para $i$ suficientemente grande, $H_o$ proviene de un $\mathbf{Z}_\ell$-haz constante torcido $\mathcal{H}_i$ sobre $S_i - D_i$ (por la «constructibilidad genérica» (0.1)) y $\mathcal{H}_i/\ell\mathcal{H}_i$ es constante.

Las descomposiciones (0.3.1) y (0.4.1) dan lugar a diagramas conmutativos:

\begin{equation}\tag{1.3.1}
\begin{tikzcd}[column sep=small, row sep=large]
\mathrm{Gal}(\bar{\eta}/\eta) \arrow[d] & \supset & I \arrow[d] & \supset & P \arrow[d] \\
\pi_1(S_i-D_i,\bar{\eta}) & \supset & I_i & \supset & P_i
\end{tikzcd}
\qquad\text{y}\qquad
\begin{tikzcd}[row sep=large]
I/P \arrow[r, dash, "\sim"] \arrow[d] & \hat{\mathbf{Z}}(1) \arrow[d, equal] \\
I_i/P_i \arrow[r, dash, "\sim"] & \hat{\mathbf{Z}}(1)
\end{tikzcd}
\end{equation}

% ESTRUCTURA: La etiqueta (1.3.1) está impresa en el margen izquierdo, a la altura del centro de la
%   figura. DOS diagramas separados en una sola línea, unidos por la palabra impresa «y» entre ellos
%   («y» es texto, no un nodo). DIAGRAMA IZQUIERDO: 6 nodos, 3 flechas y 4 símbolos de relación impresos.
%   Nodos: N1 = Gal(\bar{\eta}/\eta) (fila superior, col. 1); N2 = I (fila superior, col. 2); N3 = P
%   (fila superior, col. 3); N4 = \pi_1(S_i-D_i,\bar{\eta}) (fila inferior, col. 1); N5 = I_i (fila
%   inferior, col. 2); N6 = P_i (fila inferior, col. 3). Relaciones impresas (símbolos compuestos entre
%   columnas, no flechas): \supset entre N1 y N2; \supset entre N2 y N3; \supset entre N4 y N5;
%   \supset entre N5 y N6. Flechas: A1 N1 -> N4, vertical, hacia abajo, flecha sencilla, sin etiqueta.
%   A2 N2 -> N5, vertical, hacia abajo, flecha sencilla, sin etiqueta. A3 N3 -> N6, vertical, hacia
%   abajo, flecha sencilla, sin etiqueta. Sin ganchos, dobles puntas ni trazos discontinuos. DIAGRAMA
%   DERECHO: 4 nodos, 4 conectores. Nodos: M1 = I/P (arriba a la izquierda); M2 =
%   \hat{\mathbf{Z}}(1) (arriba a la derecha); M3 = I_i/P_i (abajo a la izquierda); M4 =
%   \hat{\mathbf{Z}}(1) (abajo a la derecha). Conectores: B1 M1 — M2, horizontal, dispuesto de
%   IZQUIERDA A DERECHA pero dibujado como una REGLA SIN PUNTA DE FLECHA, con la etiqueta \sim centrada
%   ENCIMA de la regla (isomorfismo escrito como barra con tilde). B2 M3 — M4, horizontal, la misma regla
%   sin punta, con \sim centrado encima. B3 M1 -> M3, vertical, hacia abajo, flecha sencilla, sin etiqueta.
%   B4 M2 = M4, vertical, dibujado como DOBLE BARRA VERTICAL (signo de igualdad girado), sin punta ni
%   etiqueta. Marca espuria: un trazo corto de tinta diagonal se encuentra justo a la izquierda de A1,
%   aproximadamente a dos tercios de su altura; no es una etiqueta ni forma parte del diagrama.
y la representación $H_o$ de $Gal(\bar{\eta}/\eta)$ se deduce de una acción de $\pi_1(S_i-D_i,\bar{\eta})$ sobre $H_o$, trivial sobre $H_o/\ell H_o$. El grupo \sisi{$Ker(GL(H)$ $\longrightarrow$ $GL(H/$$\ell$$H))$}{$Ker(GL(H_o)$ $\longrightarrow$ $GL(H_o/$$\ell$$H_o))$} es un pro-$\ell$-grupo; por tanto, $P_i$ actúa trivialmente, y se aplica la variante siguiente de 1.1.

\underline{Variante} 1.4.  \underline{Retomemos las notaciones de} 0.4 \underline{y sea} $\rho$ \underline{una representación continua de} $Gal(\bar{\eta}/\eta)$ \underline{en} $\sisi{Gal}{GL}(n,\mathbf{Q}_\ell)$. \underline{Supongamos que} $k(s)$ \underline{satisface} $(\ast_\ell)$ \underline{y que} $\rho(P)$ \underline{es finito}. \underline{Entonces,} \underline{la acción de} $I$ \underline{es cuasiunipotente}.


2. \underline{Ciclos evanescentes}

2.1.    Sea $S$ un rasgo estrictamente local, y retomemos las notaciones de 0.3. Por hipótesis, aquí se tiene $s = \bar{s}$. Sea $f : X \longrightarrow S$ un esquema de tipo finito sobre $S$. El formalismo de los ciclos evanescentes es una herramienta para estudiar la relación entre las cohomologías de $X_s$ y de $X_{\bar{\eta}}$, y la acción del grupo de inercia sobre la de $X_{\bar{\eta}}$, a partir de información \underline{local} (sobre $X$) acerca del morfismo $f$.

Un resultado clave de este tipo es el teorema de especialización SGA 4 XVI 2.2, según el cual, para $f$ propio y liso, $X_s$ y $X_{\bar{\eta}}$ tienen la misma cohomología.

En la teoría trascendente, si $f : X \longrightarrow D$ es un morfismo propio de un espacio analítico $X$ en el disco $D$, los ciclos evanescentes aparecen como sigue:

a) Tras reducir $D$ si es preciso, se puede suponer que $X$ es un fibrado topológico sobre el disco perforado $D^{*}$ (teorema de isotopía de Thom). Las fibras


%%%%%%%%%%  scan idx 19  |  folio 8  %%%%%%%%%%

$X_t$ ($t \in D^{*}$) son, en particular, todas isomorfas.

b) No es irrazonable representarse la fibra especial $X_s$ como deducida de «la» fibra general $X_t$ por contracción de poliedros en $X_t$: se tendría $\pi : X_t \longrightarrow X_s$. El método de los ciclos evanescentes consiste entonces en estudiar la diferencia entre las cohomologías de $X_s$ y $X_t$ mediante la sucesión espectral de Leray de $\pi$.

La teoría geométrica explicada a continuación es muy próxima a esta teoría trascendente; hay que sustituir el punto general $t \in D$ por $\bar{\eta}$, el punto geométrico genérico de $S$. Recíprocamente, se puede modelar la teoría trascendente según la teoría geométrica, sustituyendo $t$ o $\bar{\eta}$ por el recubrimiento universal $\widetilde{D}^{*}$ de $D^{*}$, y $X_t$ o $X_{\bar{\eta}}$ por $X \times_D \widetilde{D}^{*}$. Al hacerlo, se eliminan numerosos $\varepsilon$ y $\eta$, pero a veces se pierde información.

2.2. Fijemos un anillo de torsión $\Lambda$ en el que $p$ sea invertible (por ejemplo, $\Lambda = \mathbf{Z}/\ell^n\mathbf{Z}$, con $\ell$ primo con $p$). Sean los morfismos:

\begin{center}
\begin{tikzcd}[column sep=large, row sep=huge]
X_s \arrow[r, "i"] \arrow[d, "f_s"] & X \arrow[d, "f"] & & X_{\bar{\eta}} \arrow[ll, "\bar{j}"] \arrow[d, "f_{\bar{\eta}}"] \\
s \arrow[r] & S & \eta \arrow[l] & \bar{\eta} \arrow[l]
\end{tikzcd}
\end{center}
% ESTRUCTURA: Diagrama sin etiqueta (no se imprime número de ecuación). 7 nodos, 8 flechas. Dos filas;
%   la fila inferior tiene 4 columnas, y la superior tiene 3 nodos situados sobre las columnas 1, 2 y 4
%   (la columna 3 de la fila superior está vacía). Nodos: N1 = X_s (arriba, col. 1); N2 = X (arriba,
%   col. 2); N3 = X_{\bar{\eta}} (arriba, col. 4); N4 = s (abajo, col. 1); N5 = S (abajo, col. 2);
%   N6 = \eta (abajo, col. 3); N7 = \bar{\eta} (abajo, col. 4). Flechas: A1 N1 -> N2, horizontal,
%   hacia la DERECHA, flecha sencilla, etiqueta i ENCIMA del fuste. A2 N3 -> N2, horizontal, hacia la
%   IZQUIERDA, abarca dos columnas (desde la col. 4, atravesando la col. 3, hasta la col. 2), flecha
%   sencilla, etiqueta \bar{j} (j con barra superior) ENCIMA del fuste. A3 N1 -> N4, vertical, hacia
%   ABAJO, flecha sencilla, etiqueta f_s a la DERECHA del fuste. A4 N2 -> N5, vertical, hacia ABAJO,
%   flecha sencilla, etiqueta f a la DERECHA del fuste. A5 N3 -> N7, vertical, hacia ABAJO, flecha
%   sencilla, etiqueta f_{\bar{\eta}} a la DERECHA del fuste. A6 N4 -> N5, horizontal, hacia la DERECHA,
%   flecha sencilla, sin etiqueta. A7 N6 -> N5, horizontal, hacia la IZQUIERDA, flecha sencilla, sin
%   etiqueta. A8 N7 -> N6, horizontal, hacia la IZQUIERDA, flecha sencilla, sin etiqueta. No hay ganchos,
%   dobles puntas, tildes de isomorfismo ni flechas discontinuas en este diagrama.
Los haces de ciclos evanescentes $\psi_i$ sobre $X_s$ son los

\begin{equation*}\tag{2.2.1}
\psi^i \;=\; i^{*}\, R^i\, \bar{j}_{*}\, \Lambda \;.
\end{equation*}
Si $f$ es \underline{propio}, se tiene (SGA 4 XII 5.5)

\begin{equation*}\tag{2.2.2}
i^{*} \;:\; H^p(X,R^q\,\bar{j}_{*}\,\Lambda) \xrightarrow{\ \sim\ } H^p(X_s,i^{*}R^q\,j_{*}\,\Lambda) \qquad ,
\end{equation*}
y la sucesión espectral de Leray para $\bar{j}$ se reescribe

\begin{equation*}\tag{2.2.3}
H^p(X_s,\psi^q) \Longrightarrow H^{p+q}(X_{\bar{\eta}},\Lambda) \qquad .
\end{equation*}
Cuando no sucede nada en el infinito, a veces puede obtenerse esta sucesión espectral sin la hipótesis de que el morfismo sea propio. Que se necesita alguna hipótesis en el infinito se ve ya en el caso trivial $X_s = \emptyset$.

%%%%%%%%%%  scan idx 20  |  folio 9  %%%%%%%%%%

De las definiciones se deduce inmediatamente:

\underline{Proposición} 2.3.  \underline{Sean} $\bar{x}$ \underline{un punto geométrico de} $X_s$ \underline{(por ejemplo, un punto cerrado),} $X_{(\bar{x})}$ \underline{el henselizado estricto de} $X$ \underline{en} $\bar{x}$ \underline{y} $X_{(\bar{x})\bar{\eta}}$ \underline{la fibra geométrica genérica de} \underline{la proyección de} $X_{(\bar{x})}$ \underline{sobre} $S$. \underline{Se tiene}

\begin{equation*}
\psi^i \;=\; H^i(X_{(\bar{x})\bar{\eta}},\Lambda)
\end{equation*}
\underline{Corolario} 2.4.  \underline{Allí donde} $f$ \underline{es liso,} $\psi^i = 0$ \underline{para} $i > 0$ \underline{y} $\psi^{0} = \Lambda$.

Esto es SGA 4 XV 2.1.

El grupo Gal($\bar{\eta}/\eta$) (igual, por hipótesis, al grupo de inercia) actúa por transporte de estructura sobre los $\psi^i$, y

\underline{Escolio} 2.5.  \underline{La sucesión espectral} 2.2.3 \underline{es} Gal($\bar{\eta}/\eta$)-\underline{equivariante}.

\underline{Variante} 2.6.  Sean $\bar{S}$ $=$ Spec($\bar{V}$) y $\bar{X} = X \times_S \bar{S}$. $X_{\bar{\eta}}$ es el complemento de $X_s = X_{\bar{s}}$ en $\bar{X}$. Pongamos

\begin{equation*}
\phi^i \;=\; \underline{H}^{\,i+1}_{X_s}\,(\bar{X}\;,\Lambda\;)
\end{equation*}
(haces de cohomología con soporte). Se tiene una sucesión exacta de haces

\begin{equation*}
0 \longrightarrow \phi^{-1} \longrightarrow \Lambda \longrightarrow \psi^{0} \longrightarrow \phi^{0} \longrightarrow 0
\end{equation*}
y, para $i > 0$, isomorfismos $\psi^i \xrightarrow{\ \sim\ } \phi^i$. Para $f$ liso, los $\phi^i$ son nulos.

Pongamos $\phi^i_{gl} = H^{\sisi{i}{i+1}}_{X_s}(\bar{X}\;,\Lambda)$ (análogo global de los $\phi^i$). Se dispone de una sucesión espectral

\begin{equation*}\tag{2.6.1}
H^p(X_s,\phi^q) \Longrightarrow \phi^{p+q}_{gl}
\end{equation*}
y de una sucesión exacta larga de cohomología que, para $f$ propio, se reescribe

\begin{equation*}\tag{2.6.2}
\dots \longrightarrow H^i(X_s\;,\Lambda\;) \longrightarrow H^i(X_{\bar{\eta}}\;,\Lambda\;) \longrightarrow \phi^i_{gl} \longrightarrow \dots
\end{equation*}
\underline{Variante} 2.7.  Sea $K_t$ la extensión moderadamente ramificada máxima de $K$ en $\bar{K}$. Se pone $X_t = X \times_S$ Spec($K_t$). Sea $\bar{j}_t$ la proyección de $X_t$ en $X$. Se definen variantes moderadas ("tame") de los $\psi^i$ poniendo


%%%%%%%%%%  scan idx 21  |  folio 10  %%%%%%%%%%

\begin{equation*}
\psi^{i}_{t} \;=\; R^{i}\,\bar{j}_{t*}\,\Lambda
\end{equation*}
Utilizando que $P$ es un pro-$p$-grupo, con $p$ invertible en $\Lambda$, se obtiene

\begin{equation}\tag{2.7.1}
H^{*}(X_{t},\Lambda) \;=\; H^{*}(X_{\bar{\eta}},\Lambda)^{P} \qquad ,
\end{equation}
\begin{equation}\tag{2.7.2}
\psi^{i}_{t} \;=\; \psi^{iP} \qquad .
\end{equation}
Para $f$ propio, la sucesión espectral de Leray de $\bar{j}_{t}$ se deduce de 2.2.3 pasando a los $P$-invariantes:

\begin{equation}\tag{2.7.3}
H^{p}(X_{s},\psi^{q}_{t}) \;\Longrightarrow\; H^{p+q}(X_{t},\Lambda) \;=\; H^{p+q}(X_{\bar{\eta}},\Lambda)^{P}\ .
\end{equation}

3. \underline{Demostración geométrica del teorema de monodromía}

Sea $\Lambda$ como en 2.2.

\underline{Conjetura} 3.1 (\underline{pureza}). \underline{Sean} $A$ \underline{un anillo excelente estrictamente local} \sisi{(0.2)}{(0.0.2)} \underline{regular y} $D$ \underline{un divisor regular de} Spec($A$) \underline{definido por un parámetro} $x$. \underline{Sea} $U$ $=$ Spec($A$) $-$ $D$ $=$ Spec($A\left[\sisi{1/X}{1/x}\right]$). \underline{Se tiene}

\begin{equation}\tag{3.1.1}
H^{q}(U,\Lambda) \;=\; \left\{\begin{array}{ll} \Lambda & \text{para } q = 0 \\[4pt] \Lambda(-1) & \text{para } q = 1 \text{ (notación } (0.\ 0.6)) \\[4pt] 0 & \text{para } q > 1 \end{array}\right.
\end{equation}
He aquí casos en los que se conoce esta conjetura (SGA 4 XIX):

para $q = 0$ o $1$, para $A$ de característica $0$, para $A$ de igual característica $p$ cuando se dispone de la resolución de singularidades en dimensión $\leqslant$ dim($A$) y en característica $p$, o para (Spec($A$),$D$) un par liso (teorema de pureza relativa, SGA 4 XVI 3.7).

Cuando se dispone de la pureza y $D = \sum\limits_{i \in B} D_{i}$ es un divisor con cruzamientos normales en Spec $A$, definido por una parte de un sistema regular de parámetros, la cohomología de $U$ $=$ Spec($A$) $-$ $D$ viene dada por


%%%%%%%%%%  scan idx 22  |  folio 11  %%%%%%%%%%

\begin{equation}\tag{3.1.2}
\begin{aligned} H^{0}(U,\Lambda) &= \Lambda \\[4pt] H^{1}(U,\Lambda) &= \Lambda(-1)^{B} \\[4pt] H^{q}(U,\Lambda) &= \overset{q}{\bigwedge} H^{1}(U,\Lambda) \qquad \text{(dado por el producto cup)} \end{aligned}
\end{equation}
Estas fórmulas son idénticas a las que, sobre $\mathbb{C}$, dan la cohomología de $D^{*B} \times D^{k}$.

3.2. Sea $S$ un rasgo estrictamente local. Conservamos las notaciones (0.0.1) y (0.0.3). Sea $f : X \longrightarrow S$ un morfismo de tipo finito. Supongamos que $X$ es regular, que $X_{\eta}$ es liso y que $X_{s}$ es un divisor con cruzamientos normales en $X$. Supondremos que este divisor es (globalmente, y no solo localmente para la topología étale) suma de divisores regulares: $(X_{s})_{\mathrm{red}} = \sum\limits_{i \in B} D_{i}$ (el caso general puede recuperarse por localización). Denotemos por $m_{i}$ sus multiplicidades: $X_{s} = \sum\limits_{i \in B} m_{i}\, D_{i}$.

Sean $x$ un punto de $X_{s}$, $\bar{x}$ un punto geométrico de $X$ localizado en $x$ (por ejemplo, $x$ cerrado, $x = \bar{x}$), $C = \{i \mid x \in D_{i}\}$ y $K$ el complejo concentrado en grados $0$ y $-1$

\begin{equation*}
K \;:\; \mathbf{Z}^{C} \;\xrightarrow{\;\;d\;\;}\; \mathbf{Z} \;:\; d((n_{j})) \;=\; \sum\limits_{j} m_{j}\; n_{j} \qquad .
\end{equation*}
\underline{Teorema} 3.3. (\underline{utiliza la pureza}). \underline{Los grupos de ciclos evanescentes moderados} $\psi^{i}_{t\bar{x}}$ \underline{son los siguientes}:

(i) \underline{Se tiene un isomorfismo canónico de} $\Lambda$-\underline{álgebras graduadas}

\begin{equation*}
\psi^{*}_{t\ \bar{x}} \;\simeq\; \overset{*}{\bigwedge}(H_{1}(K) \otimes \Lambda(-1)) \otimes_{\Lambda} \; \psi^{0}_{t\bar{x}} \qquad .
\end{equation*}
(ii) $\psi^{0}_{t\ \bar{x}}$ \underline{es una} $\Lambda$-\underline{álgebra} $\Lambda^{E}$, \underline{para} $E$ \underline{un conjunto sobre el cual el grupo de inercia moderada} $I_{t} = I/P = \hat{\mathbf{Z}}(1)(k(s))$ \underline{actúa transitivamente}, \underline{cuyo número de elementos es el mayor divisor primo con} $p$ \underline{de} $\#$ $H_{0}(K)$.

He aquí el método de demostración.

a) Sean $X_{(\bar{x})}$ la localización estricta de $X$ en $\bar{x}$, $U = X_{(\bar{x})} - X_{(\bar{x})s}$ (el complemento de un divisor con cruzamientos normales), $t_{j}$ una ecuación local de $D_{j}$, $X_{n} = \sisi{X_{(x)}}{X_{(\bar{x})}}\left[(t_{j}^{1/n})\right]$ (para $(n,p) = 1$; es un esquema regular) y $U_{n}$ la imagen inversa de $U$ en $X_{n}$. Según 3.1.2, se tiene $H^{q}(U_{n},\Lambda) = \overset{q}{\bigwedge}(\Lambda(-1)^{C})$; para $m = \sisi{x d}{n d}$, la aplicación «imagen inversa» de $H^{q}(U_{n},\Lambda)$ en $H^{q}(U_{m},\Lambda)$ es la multiplicación


%%%%%%%%%%  scan idx 23  |  folio 12  %%%%%%%%%%

por $d^{q}$. Poniendo $\widetilde{U} = \varprojlim\limits_{(n,p)=1} U_{n}$, se tiene, pues,

\begin{equation}\tag{3.3.1}
\left\{\begin{array}{l} H^{0}(\widetilde{U},\Lambda) \;=\; \Lambda \\[8pt] H^{q}(\widetilde{U},\Lambda) \;=\; \varinjlim\limits_{n} H^{q}(U_{n},\Lambda) \;=\; 0 \ \ \text{para}\ \ q \neq 0 \ . \end{array}\right.
\end{equation}
b) $\widetilde{U}$ es un prorecubrimiento de $U$ de grupo $\hat{\mathbf{Z}}(1)(k(s))^{C}$. Por definición, se tiene $\psi^{*}_{t\ \bar{x}} = H^{*}(U_{\eta_{t}} ,\Lambda)$; $\widetilde{U}$ es un prorecubrimiento de grupo $H_{1}(K) \otimes \hat{\mathbf{Z}}(1)(k(s))$ de cada una de las componentes conexas de $U_{\eta_{t}}$, y estas son permutadas transitivamente por el grupo de inercia moderada. De (3.3.1) y de la sucesión espectral de Hochschild--Serre se deducen entonces las fórmulas 3.3.

En términos gráficos, puede decirse que $U_{\bar{\eta}}$ es la fibra de un morfismo $\delta$ de un $K(\hat{\mathbf{Z}}(1)(k(s))^{C},1)$ cohomológico $U$ en el $K(\hat{\mathbf{Z}}(1)(k(s)),1)$ cohomológico $\eta$, y que este morfismo $\delta$ induce $d = \Sigma\, m_{j}\, n_{j} \;:\; \hat{\mathbf{Z}}(1)(k(s))^{C} \longrightarrow \hat{\mathbf{Z}}(1)(k(s))$.

\underline{Corolario} 3.4. (\underline{utiliza la pureza}). \underline{Supongamos que} $f$ \underline{es propio}. \underline{Sea} $N$ \underline{el mínimo común múltiplo de las multiplicidades} $m_{j}$ (j$\in$B). \underline{Sea} $q \geqslant 0$, \underline{y} $q'$ \underline{el mínimo de} $q$ \underline{y del número máximo de divisores} $D_{j}$ \underline{que pasan por un mismo punto}. \underline{Para} $T$ \underline{en el grupo de inercia moderada} $I_{t} =I/P$, \underline{se tiene}

\begin{equation*}
(T^{N} - I)^{\sisi{q'}{q'+1}} \;=\; 0 \quad \underline{\text{sobre}} \ \ H^{q}(X_{\bar{\eta}} \,,\Lambda\ )^{P} \quad .
\end{equation*}
Se deduce inmediatamente de 3.3 y de la variante 2.7 de 2.5.

Los resultados de pureza necesarios están disponibles para los $H^{1}$, en característica $0$, o para $S$ esencialmente de tipo finito sobre un cuerpo.

\underline{Teorema} 3.5. \underline{Sea} $C$ \underline{una curva proyectiva y lisa sobre el cuerpo de fracciones} k($\eta$) \underline{de un anillo de valoración discreta} $V$. \underline{El grupo de inercia} $I$ \underline{actúa sobre} $H^{1}(C_{\bar{\eta}},\mathbf{Q}_{\ell})$ \underline{por automorfismos cuasiunipotentes de escalón} $2$: \underline{existe} $N$ \underline{tal que,} \underline{para} $T \in I$, \underline{se tenga}

\begin{equation*}
(T^{N} - I)^{2} \;=\; 0 \quad \text{sobre} \ \ H^{1}(C_{\bar{\eta}},\mathbf{Q}_{\ell}) \qquad .
\end{equation*}
Podemos reducirnos a suponer que $V$ es estrictamente local. La afirmación es que existe un subgrupo abierto $I_{1}$ de $I$ tal que $(T - I)^{2} = 0$ sobre $H^{1}$ para $T \in I_{1}$.

%%%%%%%%%%  scan idx 24  |  folio 13  %%%%%%%%%%

Podemos extender los escalares de $k(\eta)$ a una extensión finita (sustituyendo $V$ por su normalización en esta extensión). Puesto que la imagen del pro-$p$-grupo $P$ en un grupo de Lie $\ell$-ádico $GL(n,Q_\ell)$ es finita, podemos suponer, en particular, que $P$ actúa trivialmente.

Sea $C_o$ un modelo minimal de $C$ sobre $V$ (para la existencia de $C_o$, basada en el teorema de Abhyankar, véase la discusión de [3] § 2, p. 87; aquí se podría reducir al caso de $V$ excelente pasando al completado de $V$: invóquese SGA 4 XVI 1.6).

Sea $C_1$ un modelo regular de $C$, deducido de $C_o$ por explosiones, cuya fibra especial es un divisor con cruzamientos normales. Los argumentos 3.3 y 3.4 se aplican a $C_1$, $\psi_t^0$, $\psi_t^1$ y $H^1(C_{\bar\eta},\mathbb{Z}/\ell^n)^P = H^1(C_{\bar\eta},\mathbb{Z}/\ell^n)$, con $N$ independiente de $n$, y de ello se deduce 3.5.

\underline{Observación} 3.6.    Sea $A$ una variedad abeliana sobre $k(\eta)$. Puesto que $A$ es cociente de una jacobiana, se obtiene asimismo que la acción de $I$ sobre $H^1(\sisi{A_\eta}{A_{\bar\eta}},\mathbb{Q}_\ell)$ o $T_\ell(A)$ es cuasiunipotente de escalón 2.

\underline{Observación} 3.7.   La acción de $I$ sobre $H^1(X_{\bar\eta},\sisi{\mathbb{Q}}{\mathbb{Q}_\ell})$ es, de hecho, cuasiunipotente de escalón 2 para todo esquema $X$ de tipo finito sobre $\eta$.

El teorema 3.5 (o 3.6) constituye la base de la demostración del teorema de reducción estable para las variedades abelianas que se dará en la exposición IX § 3.

3.8.   Está claro que, cuando se dispone de la pureza y de la resolución de singularidades (por ejemplo, en característica 0), el método precedente demuestra el teorema de monodromía y proporciona cotas para el exponente de nilpotencia. Así, en igual característica 0, se obtiene que, para $X$ propio y liso sobre $\eta$ y $T$ perteneciente a un subgrupo abierto de $I$, se tiene

\begin{equation*}
(T - 1)^{q+1} \;=\; 0 \quad \text{sobre} \quad H^q(X_{\bar\eta}\,,\, \mathbb{Z}_\ell) \qquad ( \, T \in I_1 \subset I \, ) \; .
\end{equation*}

4. \underline{Criterios de anulación para los haces de ciclos evanescentes}

4.1.  Sean $S$ un rasgo estrictamente local y $f : X \longrightarrow S$ un esquema de tipo finito sobre $S$. Pongamos $n = \dim(X_\eta)$.

Si $X'$ es la adherencia esquemática de $X_\eta$ en $X$, se tiene:

a) Sobre $X_s - X'_s$, $\phi^{-1} = \Lambda$ y los demás $\phi^i$ son nulos;


%%%%%%%%%%  scan idx 25  |  folio 14  %%%%%%%%%%

b) sobre $X'_s$, los $\phi^i$ y $\psi^i$ relativos a $X/S$ o a $X'/S$ coinciden, y $\phi^{-1} = 0$.

Puesto que $X'$ es plano sobre $S$, esto permite a menudo reducirse al caso en que $X$ es plano sobre $S$.

\underline{Teorema} 4.2.

(i) \underline{Se tiene} $\psi^i = 0$ \underline{para} $i > n$.

(ii) \underline{Más precisamente}, \underline{para} $x$ \underline{un punto de} $X_s$ \underline{cuya adherencia sea de} \underline{dimensión} $k$ ($k = deg tr(k(x)/k(s))$) \underline{y para} $\bar{x}$ \underline{un punto geométrico localizado en} $x$, \underline{se tiene}

\begin{equation*}
\psi^i_{\bar{x}} \;=\; 0 \quad \underline{\text{para}} \quad i > n-k \;\, , \;\, \text{es decir,}
\end{equation*}
\begin{equation*}
d(\psi^i) \;\leqslant\; n-i \qquad\qquad (\underline{\text{notación de}}\ \text{SGA 4 XIV 2.1}).
\end{equation*}
Se tiene $\psi^i_{\bar{x}} = H^i(X_{(\bar{x})\bar\eta}\, ,\, \Lambda)$, y $X_{(\bar{x})\bar\eta}$ es límite proyectivo de esquemas afines de dimensión $\leqslant$ $n$ sobre $k(\bar\eta)$. La afirmación (i) se deduce, pues, del teorema de Lefschetz afín (SGA 4 XIV 3.2) por paso al límite.

Para demostrar (ii), se puede suponer que $X$ es plano sobre $S$ (4.1) y que existe un $S$-morfismo cuasifinito y plano $g : X \longrightarrow \mathbb{A}^k_S$ tal que $g(x)$ sea el punto genérico de $\mathbb{A}^k_s$. Sea $S'$ la localización estricta de $\mathbb{A}^k_S$ en $\bar{x}$. Se tiene

a) $S'_{\bar\eta}$ es el espectro de un cuerpo, y $\mathrm{Gal}(\overline{\eta'}/s'_{\bar\eta})$ es un $p$-grupo $Q$.

b) Según (i) para $g$, se tiene $H^i(X_{(x)\overline{\eta'}},\Lambda) = 0$ para $i > n-k$.

c) Se tiene $H^i(X_{(x)\bar\eta},\Lambda) = H^i(X_{(x)\overline{\eta'}},\Lambda)^Q$, de donde se deduce el teorema.

\underline{Corolario} 4.3.  \underline{Si} $f$ \underline{es propio y plano y la dimensión del lugar de no lisitud} \underline{de} $f$ \underline{en} $X_s$ \underline{es} $\leqslant d$, \underline{el morfismo de especialización} $H^i(X_s) \longrightarrow H^i(X_{\bar\eta})$ \underline{es} \underline{un isomorfismo para} $i > n+d+1$ \underline{y un epimorfismo para} $i = n+d+1$.

4.4.  Para demostrar, en ciertos casos, que los $\phi^i$ son nulos para $i$ pequeño, nos apoyaremos en la teoría de la dualidad local (directamente o mediante teoremas locales de Lefschetz de SGA 2 XIV). Esta herramienta solo está disponible en característica 0 (o en igual característica $p$ cuando se dispone de la resolución de singularidades al modo de Hironaka en las dimensiones consideradas).


%%%%%%%%%%  scan idx 26  |  folio 15  %%%%%%%%%%

\underline{Teorema} 4.5.  \underline{Supongamos que} $S$ \underline{es de característica} 0. \underline{Sean} $f : X$ $\longrightarrow$ $S$ \underline{un morfismo de tipo finito y} $x$ \underline{un punto cerrado de} $X_s$. \underline{Supongamos que}

(a) $x$ \underline{es un punto aislado de no lisitud en su fibra}

(b) $X$ \underline{es de profundidad geométrica relativa} $\geqslant$ $n$ \underline{en} $x$: \underline{en un entorno de} $x$, $X$ \underline{se identifica con un subesquema definido por} $k$ \underline{ecuaciones en un esquema liso de dimensión relativa} $N$ \underline{en} $x$, \underline{con} $n$ $\leqslant$ $N-k$.

\underline{Entonces} $\phi^i_x = 0$ \underline{para} $i < n$.

Sea $X_x$ la localización estricta de $X$ en $x$ y $V = X_x - \{x\}$. Para todo rasgo $S'$ finito sobre $S$, sean asimismo $X'_x = X_x \times_S S'$ y $V' = V \times_S S'$. Para $\sisi{S}{\bar{S}}$ la normalización de $S$ en $\bar{\eta}$, sean $\bar{X}_x = X_x \times_S \bar{S}$, $\bar{V} = V \times_S \bar{S}$. Se tiene

(1) Si $\phi^i = 0$ sobre $(X_x)_s$ (es decir, en toda generalización de $x$ en $X_s$) para $i$ $\leqslant$ $m$, entonces

\begin{equation*}
H^{i-1}_{\{x\}}(\bar{X}_x) \;=\; \widetilde{H}^{i}(\bar{V}) \;\xrightarrow{\ \sim\ }\; \widetilde{H}^{i}(\bar{V}_{\bar{\eta}}) \;=\; \phi^i \qquad\qquad \text{para}\ \ i \leqslant m\ .
\end{equation*}
Los $\phi^i$ son, en efecto, haces de cohomología con soporte en $V_s$ de $\bar{V}$; bajo las hipótesis de 4.5, \sisi{(a)}{(1)} es aplicable para $m =$ $\infty$ (SGA 4 XV 2.1).

(2) La hipótesis (b) de 4.5 es estable por cambio de base. Según SGA 2 XIV 5.6, implica que

\begin{equation*}
\widetilde{H}^{i}(V') \;=\; H^{i-1}_{\{x\}}(X'_x) \;=\; 0 \qquad\qquad \text{para}\ \ i < n\ ;
\end{equation*}
se concluye observando que $\widetilde{H}^{i}(\bar{V})$ es límite inductivo de los $\widetilde{H}^{i}(V')$.

De 4.2 y 4.5 se obtiene

\underline{Corolario} 4.6.  \underline{Para} $S$ \underline{un rasgo henseliano de característica} 0, $X/S$ \underline{una intersección completa relativa de dimensión relativa} $n$ \underline{y} $x$ $\in$ $X_s$ \underline{un punto aislado de no lisitud en su fibra}, \underline{se tiene} $\phi^i_x = 0$ \underline{para} $i$ $\neq$ $n$.

\underline{Corolario} 4.7.  \underline{Bajo las hipótesis de} 4.6, $\phi^n_x$ \underline{es un} $\Lambda$-\underline{módulo plano}.

Para $\Lambda'$ una $\Lambda$-álgebra finita, un argumento formal general (SGA 4 XVII 5.2.11) proporciona, en efecto, para los $\phi^i_{\sisi{\Lambda}{\Lambda'}}$ relativos a $\Lambda'$,


%%%%%%%%%%  scan idx 27  |  folio 16  %%%%%%%%%%

\begin{equation*}
\phi^{n-i}_{\Lambda'} \;=\; \mathrm{Tor}^{\Lambda}_{i}\ (\Lambda',\phi^n_{\Lambda})
\end{equation*}
y se aplica 4.6 a $\Lambda'$.

\underline{Variante} 4.8.  En 4.5, suprimamos la condición (a) y sea $Z$ el \sisi{lugar de no de}{lugar de no lisitud de} $f$. Todavía se puede concluir

\begin{equation*}
\phi^i_x \;=\; 0 \qquad \text{para}\ \ i < n - \dim Z_s\ .
\end{equation*}
Se procede por inducción descendente sobre dim $\overline{\{x\}}$. La hipótesis de inducción permite aplicar 4.5 (1) con $m = n -$ dim $Z_s$, y se concluye la demostración como antes.

4.9.  Es posible obtener un criterio de anulación que solo utilice información sobre las fibras geométricas de $f$. No obstante, en el caso de las intersecciones completas, el resultado es una unidad peor que 4.5.

\underline{Proposición} 4.10.  \underline{Bajo las hipótesis generales de} 4.5, \underline{supongamos que}

(a)  $x$ \underline{es un punto aislado de no lisitud en su fibra.}

(b) \underline{Para todo punto} $y$ \underline{de} $X_{\bar{\eta}}$, \underline{cuya adherencia en} $X_{\bar{\eta}}$ \underline{sea de dimensión} $d(y)$, \underline{que sea una generalización de} $x$, \underline{se tiene} \sisi{prof ét$_y(X_{\bar{\eta}})$ $\geqslant$ $n - d(g)$, es decir, $H^i_{\{y\}}(X_{\bar{\eta}}) = 0$ \underline{para} $i$ $\leqslant$ $n - d(g)$}{prof ét$_y(X_{\bar{\eta}})$ $\geqslant$ $n - d(y)$, es decir, $H^i_{\{y\}}(X_{\bar{\eta}}) = 0$ \underline{para} $i$ $\leqslant$ $n - d(y)$}

(c)  prof ét$_x(X_s)$ $\geqslant$ $n$,

\underline{Entonces,} $\phi^i_x = 0$ \underline{para} $i < n-1$.

Conservemos las notaciones de 4.5; se puede aplicar 4.5 (1) con $m =$ $\infty$. Según el teorema local de Lefschetz (SGA 2 XIV), la hipótesis (b) proporciona

\begin{center}
\begin{tikzcd}[row sep=large, column sep=huge]
H^i(V') \arrow[r, "\sim"] & H^i(V_s) & {(\ i < n\text{-}1\ )} \\
H^i(V') \arrow[r, hook] & H^i(V_s) & {(\ i = n\text{-}1\ )}
\end{tikzcd}
\end{center}
% ESTRUCTURA: Dos líneas de fórmula, dispuestas en 3 columnas (la 3.ª columna es la condición entre
%   paréntesis, alineada a la derecha, sin flecha). 6 nodos en total, 2 flechas en total. NODOS (6): N1 =
%   fila 1, col. 1: H^i(V') N2 = fila 1, col. 2: H^i(V_s) N3 = fila 1, col. 3: ( i < n-1 )
%   [anotación en redonda, entre paréntesis espaciados, NO un nodo matemático enlazado] N4 = fila 2, col.
%   1: H^i(V') N5 = fila 2, col. 2: H^i(V_s) N6 = fila 2, col. 3: ( i = n-1 ) [anotación en redonda,
%   entre paréntesis espaciados, NO un nodo matemático enlazado] FLECHAS (2): A1: N1 --> N2.
%   Horizontal, de izquierda a derecha. Trazo continuo, fuste largo, punta sencilla a la derecha.
%   Etiqueta «∼» (tilde) colocada ENCIMA del fuste, en el centro. Tipo: flecha ordinaria marcada como
%   isomorfismo por la tilde superpuesta (no es de doble punta ni un \xrightarrow{\sim} tipográfico;
%   tilde manuscrita encima). A2: N4 --> N5. Horizontal, de izquierda a derecha. Trazo continuo, fuste
%   largo, punta sencilla a la derecha y COLA EN GANCHO orientada a la izquierda (glifo «⊂» unido al
%   inicio del fuste): flecha con gancho / monomorfismo (\hookrightarrow). Sin etiqueta. NINGUNA flecha
%   vertical, diagonal, discontinua ni de doble punta. Ninguna etiqueta lateral. Las dos líneas no están
%   enlazadas entre sí por flecha alguna.
y se concluye mediante (c) que $H^i(\bar{V}) = \varinjlim H^i(V') = 0$ para $i$ $\leqslant$ $n-1$, de donde se obtiene 4.10.

4.11.  Como en 4.8, se puede dar una variante de 4.10 que no suponga (a).

%%%%%%%%%%  scan idx 28  |  folio 17  %%%%%%%%%%


5.  \underline{Acción de la monodromía sobre los} $\pi_1$

En este § (= exposición de Grothendieck del 19/11/1968), haremos uso del teorema de reducción semiestable para las curvas. Artin y Winters [2] dieron una demostración directa de este teorema, que también puede deducirse [3] del teorema análogo para las variedades abelianas (IX § 3, utilizando 3.6). Utilizaremos asimismo la teoría de Schlessinger (exposición VI de Rim).

5.1.  Sean $k$ un cuerpo algebraicamente cerrado, $C$ una curva proyectiva sobre $k$ y $x_1 \ldots x_n$ un conjunto finito de puntos cerrados de $C$. Se dice que $X = (C ; x_1,.., x_n)$ es \underline{estable} si

(a) $C$ es reducida y sus únicas singularidades son puntos dobles con tangentes distintas;

(b) los $x_i$ son distintos y son puntos lisos de $C$;

(c) si $C_1$ es una componente irreducible de $C$, y $E$ es el conjunto de puntos de la normalización $C_1'$ de $C_1$ situados sobre un punto singular de $\sisi{C_1}{C}$ o sobre uno de los $x_i$, entonces, si $C_1'$ es de género $0$ (resp. $1$), $E$ tiene al menos $3$ (resp. $1$) elementos.

La condición (c) equivale a cada una de las siguientes:

(c')  $X$ no tiene automorfismos infinitesimales no triviales.

(c'')  Si $\omega$ es el haz invertible dualizante, el haz invertible $\omega' = \omega(\Sigma\ x_i)$ es amplio.

Un $S$-esquema $C$, provisto de secciones $x_1$, ..., $x_n$, es una \underline{curva estable} sobre $S$ si sus fibras geométricas son curvas estables. En [3] §§ 1,2, solo se considera el caso $n = 0$. La extensión de los resultados de loc. cit. al caso general es inmediata:

(1) Para $X/k$ estable, $H^1(C,\omega'^{\,\otimes n}) = 0$ para $n \geqslant 2$, y $\omega'^{\,\otimes n}$ es muy amplio para $n \geqslant 3$ (véase [3] 1.2).

(2) El esquema formal $M$ de módulos de $X$ es liso, y $C$ sigue siendo singular sobre un divisor con cruzamientos normales de $M$.


%%%%%%%%%%  scan idx 29  |  folio 18  %%%%%%%%%%

(3) Para $X$ e $Y$ estables sobre $S$, $\mathrm{I}\underline{\mathrm{som}}_{S}(X,Y)$ es finito y no ramificado sobre $S$ (isomorfismos compatibles con los puntos marcados).

(4) Para $X/k$ estable, $\mathrm{Aut}_k(X) \longrightarrow \mathrm{Aut}_k\ \mathrm{Pic}^{\circ}(C)$ es inyectivo.

Se tiene, por último,

\underline{Proposición} 5.1.1.  \underline{Sea} $X_\eta = (C_\eta ; x_1, \ldots, x_n)$ \underline{estable sobre el punto genérico} $\eta$ \underline{de un rasgo} $S$, \underline{con} $C_\eta$ \underline{lisa. Existe una extensión finita separable} $\eta'$ \underline{de} $\eta$ \underline{y una curva estable} $X'$ \underline{sobre la normalización} $S'$ \underline{de} $S$ \underline{en} $\sisi{\eta}{\eta'}$ \underline{tal que}

\begin{equation*}
X' \otimes_{S'} \eta' = X_{\eta} \otimes_{\eta} \eta' \ .
\end{equation*}
Esto se deduce del teorema usual de reducción semiestable, que permite tomar $\eta'$ de modo que la fibra especial geométrica del modelo minimal $C''$ de $C_{\eta'}$ sobre $S'$ satisfaga (a). A continuación se explotan, de manera iterada si es preciso, puntos \underline{lisos} de $C''$ para que se satisfaga (b). Se contraen entonces las cadenas de componentes racionales lisas de \sisi{autointersección dos}{autointersección $-2$} que no contienen ningún $\overline{x}_i$ de la fibra especial y se obtiene $C'$ (véase [3], demostración de 2.3, p. 88).

Se puede comprobar como en loc. cit. que $C'$ existe en cuanto la jacobiana de $C_\eta$ tiene reducción estable. Por otra parte, la hipótesis de lisitud sobre $C_\eta$ es superflua;

5.2.  Sea $X_\eta$ un esquema geométricamente conexo de tipo finito sobre el cuerpo de fracciones de un rasgo estrictamente local $S$. Para $\xi$ un punto geométrico de $X_{\overline{\eta}}$ se tiene una sucesión exacta

\begin{equation}\tag{5.2.1}
\begin{tikzcd}[column sep=large, row sep=tiny]
0 \sisi{\arrow[r, dash]}{\arrow[r]} & \pi_1(X_{\overline{\eta}},\xi) \arrow[r] & \pi_1(X,\xi) \arrow[r] & \mathrm{Gal}(\overline{\eta}/\eta) \arrow[r] & 0
\end{tikzcd}
\end{equation}\ ,
% ESTRUCTURA: Una fila, 5 nodos, 4 flechas. NODOS (de izquierda a derecha): N1 = 0 N2 =
%   \pi_1(X_{\overline{\eta}},\xi) [X lleva el subíndice eta con barra; la barra se extiende a la derecha
%   de X y eta queda debajo] N3 = \pi_1(X,\xi) [impreso SIN subíndice alguno en X] N4 =
%   \mathrm{Gal}(\overline{\eta}/\eta) [barra sobre la primera eta solamente] N5 = 0 FLECHAS (todas
%   horizontales, de izquierda a derecha, sin etiqueta, continuas, ninguna con gancho / ninguna de trazos
%   discontinuos / ninguna de doble punta): A1: N1 --> N2: impresa como REGLA SIN PUNTA DE FLECHA (la
%   punta está ausente en el original mecanografiado) A2: N2 --> N3: línea continua con punta sencilla a
%   la derecha A3: N3 --> N4: línea continua con punta sencilla a la derecha A4: N4 --> N5: línea
%   continua con punta sencilla a la derecha PUNTUACIÓN: tras N5 aparece una coma, separada ampliamente a
%   la derecha del 0 final. ETIQUETA: (5.2.1) impresa al ras del margen izquierdo.
de donde se obtiene un homomorfismo de $\mathrm{Gal}(\overline{\eta}/\eta)$ en el grupo de automorfismos exteriores de $\pi_1(X_{\overline{\eta}})$. Pasando al mayor cociente primo con $p$, $\pi_1^{(p)}(X_{\overline{\eta}})$, de $\pi_1(X_{\overline{\eta}})$, se obtiene

\begin{equation}\tag{5.2.2}
\begin{tikzcd}[column sep=large, row sep=tiny]
\mathrm{Gal}(\overline{\eta}/\eta) \arrow[r] & \mathrm{Aut}\ \mathrm{ext}\ (\pi_1^{(p)}\ (X_{\overline{\eta}}))
\end{tikzcd}
\end{equation}\ .
% ESTRUCTURA: Una fila, 2 nodos, 1 flecha. NODOS (de izquierda a derecha): N1 =
%   \mathrm{Gal}(\overline{\eta}/\eta) [barra sobre la primera eta solamente] N2 = \mathrm{Aut}\
%   \mathrm{ext}\ (\pi_1^{(p)}\ (X_{\overline{\eta}})) [«Aut ext» escrito como dos palabras en redonda;
%   el superíndice (p) se sitúa sobre el subíndice 1 de \pi_1] FLECHAS: A1: N1 --> N2, horizontal, de
%   izquierda a derecha, línea continua con una punta sencilla a la derecha, sin etiqueta; no tiene
%   gancho, no es discontinua ni tiene doble punta. PUNTUACIÓN: tras N2 aparece un punto, separado por un
%   espacio. ETIQUETA: (5.2.2) impresa al ras del margen izquierdo; el cuerpo de la fórmula está sangrado
%   aproximadamente hasta la mitad de la caja.
Supongamos que $X_\eta$ tiene un punto racional $x$, y sea $\overline{x}$ el punto geométrico $\overline{\eta} \longrightarrow \eta \xrightarrow{\ x\ } X_\eta$. Para $\xi = \overline{x}$, la sucesión exacta \sisi{(5.2.2)}{(5.2.1)} se escinde canónicamente, de donde


%%%%%%%%%%  scan idx 30  |  folio 19  %%%%%%%%%%

\begin{equation}\tag{5.2.3}
[x] \ : \ \mathrm{Gal}(\bar{\eta}/\eta) \longrightarrow \mathrm{Aut}\ (\pi_1(X_{\bar{\eta}},\bar{x})) \longrightarrow \mathrm{Aut}\ (\pi_1^{(p)}(X_{\bar{\eta}},\bar{x}))\ .
\end{equation}
\underline{Teorema} 5.3. \underline{La imagen de} $P$ \underline{por} 5.2.2 (\underline{o} 5.2.3, \underline{lo que es equivalente}) \underline{es finita}.


5.3.1. \underline{Reducción al caso en que} $X_\eta$ \underline{es geométricamente normal y entero.}

Sean $X'_\eta$ la normalización de $X_\eta$, $(X'_{i\eta}\,)_i$ la familia de componentes conexas de $X'_\eta$, $X''_\eta = X'_\eta \times_{X_\eta} X'_\eta$ y $(X''_{j\eta}\,)_j$ las componentes conexas de $X''_\eta$. Tras pasar, si es preciso, a una extensión finita de $k(\eta)$, se puede suponer que los $X'_{i\eta}$ son geométricamente normales y enteros, con un punto racional $x_i$, y que los $X''_{j\eta}$ son geométricamente conexos, con un punto racional $y_j$. Para $k = 1$ o $2$ y $p_k(y_j) \in X'_{i\eta}$, escojamos una clase de caminos $\ell_{kij}$ de $\mathrm{pr}_k(\bar{y}_j)$ a $\bar{x}_i$: es un elemento de un torsor bajo $\pi_1(X'_{i\bar{\eta}},\bar{x}_i)$. Sea $\bar{\ell}_{kij}$ el camino «primo con $p$» correspondiente (el mismo módulo $\mathrm{Ker}(\pi \longrightarrow \pi^{(p)})$\sisi{}{)}. Se puede tomar $\bar{\ell}_{kij}$ invariante por $P$: el conjunto de caminos primos con $p$ es $\varprojlim$ de conjuntos finitos de orden primo con $p$ sobre los que actúa el pro-$p$-grupo $P$.

Por consiguiente, si un subgrupo de índice finito $P'$ de $P$ actúa trivialmente sobre los $\pi_1^{(p)}(X'_{i\bar{\eta}},\bar{x}_i)$, se deduce de las fórmulas de Van Kampen (es decir, del hecho de que $\sisi{X'_{i\bar{\eta}} \longrightarrow X_{i\bar{\eta}}}{X'_{\bar{\eta}} \longrightarrow X_{\bar{\eta}}}$ es de descenso efectivo para los recubrimientos étale) que $P'$ actúa trivialmente sobre $\pi_1^{(p)}(X_{\bar{\eta}},\bar{x})$.

5.3.2. Para $X_{\bar{\eta}}$ normal y $U$ un abierto, $\pi_1(U)$ se envía sobre $\pi_1(X_{\bar{\eta}})$; se puede, pues, suponer que $X_{\bar{\eta}}$ es liso y afín. Cortando por secciones hiperplanas \underline{genéricas} (esto sustituye $S$ por la localización estricta de $S$ en el punto genérico de la fibra especial de un $\mathbb{A}^n_S$) y aplicando Bertini, se reduce al caso en que $X_\eta$ es \underline{liso, geométricamente conexo y de dimensión uno}. Tras pasar, si es preciso, a una extensión finita de $k(\eta)$, se puede suponer que $X_\eta$ es el complemento, en una curva proyectiva y lisa $\bar{X}_\eta$, de un conjunto finito de puntos racionales $x_i$ (cuyo número podemos aumentar a voluntad\sisi{}{)}. Aplicando el teorema de reducción semiestable (5.1.1), se reduce al caso en que existe una curva estable $X = (C ; x_1 ,\ldots, x_n)$ sobre $S$ tal que $X_\eta = C_\eta - \{x_1 ,\ldots, x_n\}$. Se aplica entonces el resultado siguiente


%%%%%%%%%%  scan idx 31  |  folio 20  %%%%%%%%%%

\underline{Teorema} 5.4. \underline{Sean} $S$ \underline{un esquema estrictamente local,} $f : \bar{X} \longrightarrow S$ \underline{un morfismo propio y plano, cuya fibra especial es una curva cuyos únicos puntos de no lisitud son puntos dobles ordinarios, y} $Y$ \underline{un subesquema de} $\bar{X}$, \underline{reunión de un número finito de secciones disjuntas contenidas en el abierto de lisitud}. Sea $x$ \underline{un punto racional de} $X = \bar{X} - Y$. \underline{Sea} $U$ \underline{el abierto de} $S$ \underline{sobre el cual} $X$ \underline{es liso, y} $\xi$ \underline{un punto geométrico de} $U$. \underline{Entonces, la imagen de} $\pi_1(U,\xi)$ \underline{en} $\mathrm{Aut}(\pi_1^{(p)}(X_\xi,\ x\xi))$ \underline{es abeliana y de orden primo con} $p$.

Se observará que $\pi_1(U,\xi)$ solo puede actuar sobre $\pi_1^{(p)}(X_\xi,x\xi)$ porque los $\pi_1^{(p)}$ satisfacen el teorema de especialización.

Se reduce al caso en que $S$ es completo, y después al caso universal en que $S$ es un esquema formal de módulos para la fibra especial, provista de sus secciones. $S$ es entonces un anillo de series formales sobre un anillo de Cohen: $S = \mathrm{Spec}(W[[t_1 \ldots t_n]]\ )$, y $U$ es el complemento de un divisor con cruzamientos normales de ecuación $t_1 \ldots t_m = 0$. En este caso universal, según el lema de Abhyankar, $\pi_1(U,\xi)$ es abeliano y de orden primo con $p$, de donde se deduce el teorema.

\underline{Observación} 5.5. Se puede demostrar un resultado análogo a 5.4 en dimensión relativa arbitraria considerando una sección plana genérica de dimensión uno y aplicando Bertini.


6. \underline{Apéndice}: \underline{demostración aritmética del teorema de reducción estable}

(por P. Deligne)

Sean $S$ un rasgo y $A_\eta$ una variedad abeliana sobre $k(\eta)$, de dimensión $d \neq 0$. Para $S' \longrightarrow S$ un morfismo local de rasgos, se denota por $A_{\eta'}$ la variedad abeliana sobre $k(\eta')$ deducida de $A_\eta$ por extensión de escalares. Sean $\mathcal{A}_{S'}$ el modelo de Néron de $A_{\eta'}$ sobre $S'$, \sisi{$\mathcal{A}_{S's'}$ su fibra especial y $\mathcal{A}^{\circ}_{S's'}$ la componente neutra de $\mathcal{A}_{S's'}$}{$\mathcal{A}_{S',s'}$ su fibra especial y $\mathcal{A}^{\circ}_{S',s'}$ la componente neutra de $\mathcal{A}_{S',s'}$}. El teorema de reducción estable es el siguiente.

%%%%%%%%%%  scan idx 32  |  folio 21  %%%%%%%%%%

\underline{Teorema} 6.1. \underline{Para} $S'$ \underline{adecuado,} \sisi{$A_{\eta'}$ \underline{tiene reducción estable,} \underline{es decir,} $\mathcal{A}^{\circ}_{S's'}$}{$A_{\eta'}$ \underline{tiene reducción estable,} \underline{es decir,} $\mathcal{A}^{\circ}_{S',s'}$} \underline{es una extensión de una variedad abeliana por un toro.}

Se deduce con bastante facilidad de 6.1 que puede tomarse $S'$ de modo que $k(\eta')$ sea una extensión separable finita de $k(\eta)$. Supongamos primero que el cuerpo residual de $S$ es de tipo finito sobre el cuerpo primo (el mismo argumento funcionaría para $k(s)$ puramente inseparable sobre un cuerpo de esta clase).

Sean $\ell$ un número primo distinto del exponente característico residual $p$, $T_{\ell}(A)$ el módulo de Tate, $V_{\ell}(A) = T_{\ell}(A) \otimes \mathbf{Q}_{\ell}$ y $\rho$ : $\mathrm{Gal}(\bar{\eta}/\eta) \longrightarrow \mathrm{GL}(T_{\ell}(A))$ la representación natural. Una extensión preliminar de los escalares nos reduce al caso en que el grupo de inercia $I$ actúa de manera unipotente (1.2), por tanto a través de su mayor cociente pro-$\ell$-grupo $\mathbf{Z}_{\ell}(1)$ (0.3). Sean $T$ la imagen en $\mathrm{GL}(T_{\ell}(A))$ de un generador de $\mathbf{Z}_{\ell}(1)$ y $r$ el mayor entero $\geqslant 0$ tal que $(T-1)^{r} \neq 0$. Se tiene:

\underline{Lema} 6.2. \underline{La aplicación}

\begin{equation*}
\sigma \in I \ , \ v \in V_{\ell}(A) \longmapsto (\rho(\sigma)-1)^{r} \ (v)
\end{equation*}
\underline{induce los morfismos y el isomorfismo}

\begin{equation}\tag{6.2.1}
\begin{tikzcd}[column sep=huge, row sep=huge]
V_{\ell}(A)_{I} \otimes \mathbf{Q}_{\ell}(r) \arrow[r, "M"] \arrow[d] & V_{\ell}(A)^{I} \\
\sisi{V_{\ell}(A)/\mathrm{Ker}(T-1)^{r} \otimes \mathbf{Q}_{\ell}(r)}{\bigl(V_{\ell}(A)/\mathrm{Ker}(T-1)^{r}\bigr) \otimes \mathbf{Q}_{\ell}(r)} \arrow[r, "\sim"] & \mathrm{Im}(T-1)^{r} \arrow[u]
\end{tikzcd}\quad .
\end{equation}
% ESTRUCTURA: 4 nodos, 4 flechas. La etiqueta «(6.2.1)» está impresa al ras del margen izquierdo,
%   centrada verticalmente junto al diagrama; se imprime un punto «.» en la línea de base de la fila
%   inferior, bien a la derecha del nodo inferior derecho. NODOS: N1 (arriba a la izquierda):
%   $V_{\ell}(A)_{I} \otimes \mathbf{Q}_{\ell}(r)$; el subíndice I está situado bajo y a la derecha de
%   $V_{\ell}(A)$; el glifo \otimes es una cruz en un círculo sobreimpresa/rellena. N2 (arriba a la
%   derecha): $V_{\ell}(A)^{I}$ N3 (abajo a la izquierda): $V_{\ell}(A)/\mathrm{Ker}(T-1)^{r}
%   \otimes \mathbf{Q}_{\ell}(r)$ N4 (abajo a la derecha): $\mathrm{Im}(T-1)^{r}$ FLECHAS: A1: N1 -->
%   N2. Horizontal, hacia la DERECHA. Flecha sencilla continua (regla mecanografiada larga con una punta
%   a la derecha). Etiqueta «M» ENCIMA de la flecha, aproximadamente centrada. Sin gancho, sin marca de
%   mono/epimorfismo, no discontinua. A2: N1 --> N3. Vertical, hacia ABAJO (punta en el extremo inferior,
%   justo encima de N3). El fuste desciende desde debajo de \otimes en N1. Flecha sencilla continua. SIN
%   ETIQUETA. No discontinua, sin decoración. A3: N3 --> N4. Horizontal, hacia la DERECHA. Flecha sencilla
%   continua. Etiqueta «$\sim$» (una tilde, que indica un isomorfismo) ENCIMA de la flecha, a la derecha
%   del centro. No es \twoheadrightarrow, no tiene gancho ni es discontinua. A4: N4 --> N2. Vertical,
%   hacia ARRIBA (punta en el extremo superior, justo debajo de N2). Flecha sencilla continua. SIN
%   ETIQUETA. No discontinua, sin decoración.
$V_{\ell}(A)_{I}$ \underline{y} $V_{\ell}(A)^{I}$ \underline{son} $\mathrm{Gal}(\bar{s}/s)$\underline{-módulos, y} $M$ \underline{es galoisiano.}

6.3. Se dice que una representación $\ell$-ádica $\tau$ : $\mathrm{Gal}(\bar{s}/s) \longrightarrow \mathrm{GL}(V)$ es de peso $n$ ($n \in \mathbf{Z}$) si se cumple la condición siguiente:

(*) Para un modelo adecuado $X$ de $k(s)$ ($X$ = una variedad irreducible de tipo finito sobre el cuerpo primo, con $k(s) = k(X)$), $\tau$ se factoriza por $\pi_{1}(X,\bar{s})$ y, para $x$ un punto cerrado de $X$, los valores propios de $\tau(F_{x})$ son números algebraicos cuyos conjugados complejos tienen todos valor absoluto $|k(x)|^{n/2}$ ($F_{x}$ denota el Frobenius geométrico $\phi_{x}^{-1}$).


%%%%%%%%%%  scan idx 33  |  folio 22  %%%%%%%%%%

Si $V$ (resp. $W$) es de peso $n$ (resp. $m$) y $n \neq m$, se tiene $\mathrm{Hom}_{\mathrm{Gal}}(V,W) = 0$.

6.4. Se deduce de la propiedad universal del modelo de Néron que $V_{\ell}(A)^{I} = V_{\ell}(\mathcal{A}^{\circ}_{S,s})$. Si $a$ es la dimensión del mayor cociente abeliano de $\mathcal{A}^{\circ}_{S,s}$ y $m$ la dimensión de la parte multiplicativa de $\mathcal{A}^{\circ}_{S,s}$, según Weil, $V_{\ell}(A)^{I}$ es, por tanto, extensión de una representación de dimensión $2a$ y peso $-1$ por una representación de dimensión $m$ y peso $-2$.

6.5. Sea $A^{*}$ la variedad abeliana dual de $A$. Recordemos que $V_{\ell}(A^{*})$ y $V_{\ell}(A)$ (por tanto, $V_{\ell}(A^{*})^{I}$ y $V_{\ell}(A)_{I}$) están en dualidad con valores en $\mathbf{Q}_{\ell}(1)$. Si $a^{*}$ y $m^{*}$ se definen para $A^{*}$ como $a$ y $m$ para $A$, se deduce de 6.4 que $V_{\ell}(A)_{I}$ es extensión de una representación de dimensión $m^{*}$ y de peso 0 por una representación de dimensión $2a^{*}$ y de peso $-1$.

6.6. El módulo galoisiano $\mathbf{Q}_{\ell}(r)$ es de peso $-2r$. Puesto que $\mathrm{Im}(T-1)^{r} \neq 0$, se deduce de 6.4, 6.5 y del isomorfismo 6.2 que $r \leqslant 1$. Si $r = 0$, se tiene

\begin{equation*}
\left\{
\begin{array}{l}
a \ + \ m \ \leqslant \ d \\[2ex]
2a + \ m \ = \ 2d \ ,
\end{array}
\right.
\end{equation*}
por tanto, $m = 0$, $a = d$ y $A_{\eta}$ tiene buena reducción. Supongamos que $r = 1$. Se deduce entonces de 6.2 que $\mathrm{Im}(T-1)$ es de peso $-2$ y que $V_{\ell}(A)/V_{\ell}(A)^{I}$ es de peso 0; estos espacios tienen la misma dimensión $m_{1} \leqslant m$.

Se tiene entonces

\begin{equation*}
\dim V_{\ell}(A)^{I} \ = \ 2d \ - \ m_{1} \ = \ 2a \ + \ m \ ,
\end{equation*}
\begin{equation*}
2 \dim \mathcal{A}_{S,s} \ \geqslant \ 2 \ (a + m) \ \geqslant \ 2 \ a \ + \ m \ + \ m_{1} \ = \ 2 \ d \ = \ 2 \dim \mathcal{A}_{S,s} \ ,
\end{equation*}
de modo que $A_{\eta}$ tiene reducción estable ($\dim \mathcal{A}_{S,s} = a + m$). Se ha obtenido al mismo tiempo que $m_{1} = m$, es decir, que $\mathrm{Im}(T-1) \subset V_{\ell}(A)^{I} \simeq V_{\ell}(\mathcal{A}^{\circ}_{S,s})$ es el $V_{\ell}$ de la parte multiplicativa de $\mathcal{A}^{\circ}_{S,s}$.

6.7. Para tratar el caso general, utilizaremos (0.5) (véase 1.3). Nos reducimos a suponer que


%%%%%%%%%%  scan idx 34  |  folio 23  %%%%%%%%%%

(a) se cumplen las condiciones de (0.5.2) o (0.5.3);

(b) $\mathrm{Gal}(\bar{\eta}/\eta)$ actúa trivialmente sobre $A_\ell$;

(c) la acción de $I$ sobre $T_\ell(A)$ es unipotente.

Bajo estas hipótesis, demostraremos que $A_\eta$ tiene reducción estable. La hipótesis (a) permite aplicar (0.5.1) para tener $S = \varprojlim S_i$ como en 1.3, cuyas notaciones retomamos. Sea $s_i$ el punto cerrado de $S_i$.

Para $i$ suficientemente grande, la componente neutra $\mathcal{A}^\circ$ del modelo de Néron $\mathcal{A}_S$ proviene de un esquema en grupos liso con fibras conexas $\mathcal{A}_i$ sobre $S_i$, y $(\mathcal{A}_i)_\ell$ es constante sobre $S_i - \sisi{D}{D_i}$. El grupo $\pi_1(S_i - D_i , \bar{\eta})$ actúa sobre $T_\ell(A)$, y actúa trivialmente sobre $T_\ell(A)/\ell\, T_\ell(A) \simeq A_\ell$. Como en 1.3, se ve que $P_i$ actúa trivialmente. Se tiene entonces, por (1.3.1):

\begin{equation*}
T_\ell(A)_I \;=\; T_\ell(A)_{I_i} \;,\qquad T_\ell(\mathcal{A}_{i,s_i}) \;=\; T_\ell(\mathcal{A}^\circ_s) \;=\; T_\ell(A)^I \;=\; \sisi{T_\ell(A^{I_i})}{T_\ell(A)^{I_i}}
\end{equation*}
y (6.2.1) es un diagrama de $\mathrm{Gal}(\bar{s}_i/s_i)$-módulos al que se pueden aplicar los argumentos 6.4 a 6.6.


%%%%%%%%%%  scan idx 35  |  folio 24  %%%%%%%%%%


\underline{BIBLIOGRAFÍA}

[1]	M. ARTIN	Algebraic approximation of structures over complete local rings. Publ. Math. I.H.E.S. \underline{36} -1969. pp. 23-58.

[2]	M. ARTIN and G. WINTERS. Degenerate fibres and reduction of curves.

[3]	P. DELIGNE and D. MUMFORD. The \sisi{irreductibility}{irreducibility} of the space of curves of a given genus. Publ. Math. I.H.E.S. \underline{36} - 1969. pp. 75 - 110.

[4]	J.P. SERRE	Corps locaux. Publ. Math. univ. Nancago-Hermann 1962.

[5]	J.P. SERRE and J. TATE. Good reduction of abelian varieties. Ann. of Math. \underline{88} 3 - 1968 - pp. 492 - 517.
